\section{Intuición comprensiva guiada durante la inferencia}

\subsection{Visión geométrica de las trayectorias curvas}

Entendido desde la intuición geométrica, la esencia de la guía durante la inferencia es el \textbf{ruido que se desvía en una trayectoria curva}:

\begin{enumerate}
    \item Sin guía, partiendo de $x_T$ (ruido puro), el modelo difusivo sigue la dirección indicada por las puntuaciones marginales $\nabla_{x_t} \log p(x_t)$, dirigiendo la trayectoria hacia algún punto $x_0$ en el subespacio de datos.
    \item Al agregar guía, la puntuación de verosimilitud $\nabla_{x_t} \log p(y|x_t)$ aplica una fuerza adicional a la trayectoria en cada paso, haciéndola inclinarse hacia aquellos puntos de datos que satisfacen la condición $\mathcal{A}(x_0) \approx y$.
    \item Finalmente, el $x_0$ al que llega la trayectoria es tanto un punto razonable en el subespacio de datos (garantizado por las puntuaciones marginales) como uno que cumple con las restricciones condicionales (guiado por las puntuaciones de verosimilitud).
\end{enumerate}

\begin{knowledgebox}{¿Por qué los modelos difusivos son especialmente adecuados para la guía durante la inferencia?}
Los GAN y VAE son modelos de generación en un solo paso: el mapeo de ruido a datos es una única propagación hacia adelante. Una vez que este mapeo se fija, resulta difícil aplicar controles adicionales durante la inferencia.

En cambio, el proceso generativo de los modelos difusivos/flujo consiste en una serie de \textbf{pequeñas iteraciones}, donde cada paso puede ser ajustado finamente. Esto es similar a la navegación: si llegas al destino en un solo paso, solo puedes ir en línea recta; pero si lo divides en muchos pasos, puedes ajustar la dirección en cada uno según nueva información —esa es precisamente la esencia de la guía durante la inferencia.
\end{knowledgebox}

\subsection{Resumen de la sección}

La esencia geométrica de la guía durante la inferencia consiste en aplicar una fuerza correctora direccional adicional en cada paso de la trayectoria de desruido, haciendo que el resultado generado final satisfaga simultáneamente los objetivos de ser ``realista'' y ``cumplir con las condiciones''. Las características iterativas del modelo difusivo lo hacen naturalmente adecuado para este tipo de operación.
